\documentclass[11pt]{article}
\usepackage{amsmath,amssymb,amsthm,hyperref}
\newtheorem{theorem}{Theorem}
\title{The four fusion properties of signed subsets}
\author{Research certificate}
\date{October 25, 2026}
\begin{document}
\maketitle
\begin{abstract}
We prove the four algebraic assertions in Conjecture 3.2 of Gatica and
Bustamante for an arbitrary ground set, using their explicit fusion operation.
A complete Lean certificate accompanies the proof. We also distinguish these
four assertions from the stronger quasigroup division property of a loop.
\end{abstract}
Let $A$ be any set. Represent a signed subset by a pair $(P,N)$ of disjoint
subsets of $A$, recording its positive and negative atoms. For $X=(P,N)$ and
$Y=(Q,M)$, set
\[
 X\oplus Y=((P\setminus M)\cup(Q\setminus N),
 (N\setminus Q)\cup(M\setminus P)).
\]
This is Definitions 1.4, 1.8 and 1.12 of \cite{GB} in signed coordinates:
remove atoms whose opposites occur in the other input and take the union.
The source's generated space consists exactly of these disjoint pairs.
Indeed, fusion of a pure positive $P$ and pure negative $Q$ gives
$(P\setminus Q,Q\setminus P)$; conversely, a disjoint pair is obtained from
its own positive and negative subsets. The certificate proves these facts
as \texttt{generated\_is\_fusion} and \texttt{generated\_space\_exact}.
\begin{theorem}
For every ground set $A$, fusion is closed and commutative on signed subsets.
It has the unique two-sided identity $(\varnothing,\varnothing)$, and every
$(P,N)$ has the unique two-sided annihilating inverse $(N,P)$.
\end{theorem}
\begin{proof}
For closure, an atom in both output coordinates would, in the two unmixed
cases, contradict $P\cap N=\varnothing$ or $Q\cap M=\varnothing$; in the two
mixed cases it would contradict one of the set differences. Commutativity
follows by interchanging the terms of the two unions. The empty pair is an
identity directly from the formula. Any identity $e$ must satisfy
$e=0\oplus e=0$.

The proposed inverse annihilates $(P,N)$ because each output coordinate is
a union of differences of a set from itself. If $(P,N)\oplus(Q,M)=0$, the
positive coordinate gives $P\subseteq M$ and $Q\subseteq N$, while the
negative coordinate gives $N\subseteq Q$ and $M\subseteq P$. Thus $Q=N$
and $M=P$, proving uniqueness. Commutativity supplies both sides.
\end{proof}
This establishes precisely the four assertions of \cite[Conjecture 3.2]{GB}.
Taking the ground type to be the subtype of any $A\subseteq\mathbb N$
recovers its stated range; finiteness is unnecessary.

The operation also satisfies $X\oplus X=X$. For $X\ne0$, translation by $X$
sends both $X$ and $0$ to $X$, so it is not injective. It therefore lacks the
quasigroup division property of an ordinary loop. Further,
$((-X)\oplus X)\oplus X=X$ but $(-X)\oplus(X\oplus X)=0$, showing
nonassociativity. These two claims also have Lean proofs.

\paragraph{Verification and scope.}
The full four-part result is \texttt{SigmaFusion.conjecture\_3\_2} in
\texttt{formal/SigmaFusion.lean}. Its build checks every local declaration
and permits only the standard Lean/mathlib axioms. This is a certificate
for the explicit signed-subset algebra, not a foundational model of all
nonstandard class-theoretic axioms discussed in \cite{GB}.
The source and its Qeios revision retain the conjecture. The retrieved
MathDB 366458 record has no solution. We make no worldwide-priority claim.
MathDB 318649 is a related duplicate and is not counted as another outcome.
\begin{thebibliography}{9}
\bibitem{GB} I. Gatica and A. Bustamante, \emph{$\sigma$-Sets and
$\sigma$-Antisets}, arXiv:2501.00047v1, Definitions 1.2, 1.4, 1.8, 1.12;
Theorem 1.6; Conjecture 3.2.
\url{https://arxiv.org/html/2501.00047v1}.
\end{thebibliography}
\end{document}
